\scanpage{120}

% manuscript lines / position: 1–5
Inversement, \inserted{dans} une telle situation twistée $A$ est une algèbre étale quadratique sur chacun des $A_i$, donc on a une situation biquadratique \inserted{étale}, et de plus les hypothèses du début (en particulier $A^{\unsure{\gamma|\sigma|\Gamma}}=k.1_A$) sont satisfaites.

% manuscript lines / position: 6–10
\deleted{Il faut} Seule la première partie, directe, demande un peu de soin. OPS $X_1,X_2,X_3$ rev. triviaux, $I_{1X},I_{2X},I_{3X}$, et \deleted{la section} \inserted{la section construite plus haut de $X_1\wedge X_2\wedge X_3=(I_1\wedge I_2\wedge I_3)_X$} constante. \deleted{\unsure{on}} On aura alors

% manuscript lines / position: 11
\[
 \unsure{\mathscr X|\mathcal X}\subset J_X
\]

% manuscript lines / position: 12
et tout résulte des lemmes suivants

% manuscript lines / position: 13–15
\uline{Lemme} Soit $\unsure{\gamma|\mathcal G|\Gamma}$ un groupe, opérant de façon transitive sur un ens. $J$, et considérons un sous-schéma $\mathscr X$ de $X_J$ \inserted{(où $X=\operatorname{Spec} k$)} qui soit tel que

% manuscript lines / position: 16
a) $\mathscr X$ est stable sous $\gamma$, et

% manuscript lines / position: 17–18
b) le composé $k\longrightarrow\Gamma(X_J,\unsure{\underline O|\mathcal O}_{X_J})\;(\simeq k^J)$ est injectif.

% manuscript lines / position: 19–20
Alors $\mathscr X=X_J$, donc $\mathscr X$ est étale sur $X$ (et fini étale si $J$ est fini).

% manuscript lines / position: 21–24
En effet, soit $j_0\in J$, considérons la composante $X_{\unsure{0|j_0|j}}$ de $X_J$, $X_j\xrightarrow{\sim}X$, et $\mathscr X_j=X\cap\mathscr X$, il correspond à un sous-schéma fermé de $X$ donc à un idéal $\unsure{\mathfrak a|\mathcal A|\mathfrak J_j}$ de $k$. A cause de l'\unsure{homog|action} de

% manuscript lines / position: bottom left
% archive sheet number: 119
